\section{合卷：家门、资源与国家——从春秋卿族到战国政体}

鲁哀公十四年的经文在“西狩获麟”处收住；同一年，齐国陈恒杀齐简公。把这两件事并写在年表上，最容易误成一个时代忽然换了主人。其实，\eventref{SA-0481-LIN}{春秋·获麟与陈氏政变}所照见的是一段很长的积累来到显处：国君仍有名号，能够调粮、守邑、带兵和安置家臣的人，却未必仍由国君支配。春秋末叶的家族政治不是国家之外的私事；它正发生在征发、赏赐、战争和地方治理的接口上。

\subsection{一份职位怎样变成一座家门的资源}

春秋的君主不能独自穿过山河征兵，也不能亲自照管每一座城邑。卿大夫受命统军、办理外交、经营封邑，本来就是国家能运转的条件。问题不在于有无贵族，而在于这些职务所连带的田土、从属、武装和婚姻关系能否在同一家门中延续。若能延续，国君每逢战争或危机就越要借助它；家门也越能把一次公共委任留下的资源变成下一次谈判的筹码。

这不是一条处处同样的道路。齐在\eventref{SA-0548-CUIZHU}{春秋·崔杼杀齐庄公}所显示的，是卿族已经可以令公室承受内乱；到\eventref{SA-0481-LIN}{春秋·获麟与陈氏政变}，陈、田氏又越过齐君改立新君。两次行动之间不能被抹成田氏预先写好的剧本：崔、庆的败亡、田氏如何结纳城邑与人众，以及每一步权力移转的具体机制，材料并不足以逐项复原。但它们共同提示，公室一旦不能垄断军政资源，掌握资源的家族便可能在危机中决定谁保有公室名分。

直到\eventref{WQ-0386-TIANQI}{战国·田氏代齐}，周廷才承认田氏为齐侯。册命不是家族力量的起点，而是既成力量获得列国可辨名号的一环。临淄的城邑、田土与通向海滨的道路并没有因姜、田易姓而消失；变化在于，谁能以“齐”的名义调动它们。田氏的结果提醒人们，所谓“夺国”既是家门压过公室，也是把原有国的资源、风险和对外责任一并接到自己手中。

\subsection{鲁没有易姓，却把权力推得更深}

鲁的情形更能防止我们把家族政治只理解为改朝换代。三桓长期居于邑、兵和朝政之间，鲁昭公试图反制而败，\eventref{SA-0517-LU-ZHAO-EXILE}{春秋·鲁昭公出奔}遂使公室受制的结果公开化。这里的直接后果是君主失去国内立足处；中期影响却不是季氏便可无代价地统治一切。公室、三桓与各自家臣仍须争夺城邑、部众和相互承认，鲁国的名号也仍为各方所用。

因此，\eventref{SA-0502-YANGHU}{春秋·阳虎出奔}尤其值得停下来细看。阳虎是季氏家臣，他的兴败说明权力并未只从国君平直地流向三桓；掌兵掌邑的从属者也能借主家的资源越级行动。家门的扩张会制造新的依附链条，也会在链条内部产生新的竞争者。鲁没有变成一个由阳虎建立的新国，正表明资源的实际占有、礼法名分和持久统治并不是同一件事。

\subsection{晋被分开，楚则在重组前遇到阻力}

晋提供了另一种结局。它在长期霸业中让卿族分担统兵、封邑和外交，获得了迅速动员联盟的能力；\eventref{SA-0546-MIBING}{春秋·弭兵之会}以后，外部决战暂缓，这些可继承的资源又有更多时间在各家门内沉积。把弭兵直接说成分晋的原因并不妥当；它只是降低了共同作战压力的一个使能条件，晋君又缺少把土地、兵员和职位稳定收回公室的办法，才使家族竞争逐渐能够决定国家的形状。

前453年的\eventref{WQ-0453-ZHI}{战国·三家灭智}把这种竞争推到晋阳城下。韩、赵、魏合力消灭智氏后，分走的不是抽象的“权力”，而是可以继续征发的人口、土地和军队。前403年\eventref{WQ-0403-THREEJIN}{战国·三家分晋}的周王室册命，因而承认了三个事实上已能各自行动的政治体。齐是一个家门接过一国，晋则由几个家门拆开一国；两者都说明旧名分往往落在资源重组之后，而非走在前面创造资源。

楚的难题又不同于齐、晋。它拥有广阔腹地，却也须借旧贵族和地方网络跨越漫长的治理距离。战国早期的\eventref{WQ-0381-WUQI-CHU}{战国·吴起变法}试图从世卿贵族手中重新抽取部分俸禄、职位和军政资源；楚悼王死后，吴起被杀，重组便遭到反弹。不能由此断言楚的贵族从未服从中央，也不能把吴起的死亡写成所有改革必败的寓言。它所暴露的条件更具体：若改革只靠一位君主的庇护，继承时刻便会让受损家门有机会重夺资源。楚没有被一场反弹分裂，却要长期承受中心、贵族与地方之间难以一举收束的成本。

\begin{causalchain}[从家门积累到国家变形]
结构性原因是春秋国家必须通过封邑、卿族和家臣来完成征发、守城与外交；触发因素常是继承危机、君臣冲突或战争后的赏赐分配。田氏越过姜齐、三家分晋、鲁君出奔和楚的贵族反弹，是在不同政治体中出现的不同结果，不能互相充当单一原因。它们的共同中期影响，是土地、兵员与官职越来越成为必须重新界定归属的资源；战国各国由此更迫切地寻找能把这些资源接到国君和官吏体系上的办法。
\end{causalchain}

\subsection{从“谁家有兵”到“国家怎样有兵”}

这条线索不该导向一种简单的进步叙事。家族网络能提供熟人间的信任、地方知识和应急的部伍；鲁、晋、齐和楚之所以曾能运转，也正因为它们不能没有这些网络。代价是公共事务的收益和成本容易沉淀在少数家门，国君的命令又常要经过它们才能到达城邑与民众。改由君主和官吏直接组织资源，未必自动更公正，却能在列国竞争中改变动员的速度和尺度。

战国诸国所作的尝试并不相同。楚的阻力说明，触动旧贵族并不等于立刻获得稳定的中央能力；齐田氏的成功也不等于接手者已经解决国家的安全与治理。魏、韩、秦后来围绕军政、户籍、土地和官吏所作的重组，正是在这片压力中展开。尤其是\eventref{WQ-0356-SHANGYANG}{战国·商鞅变法}，并非凭一纸法令使家族关系消失，而是试图以军功、编户和县邑改变资源通向权力中心的路径；它的代价，则会更直接地落到服役、纳税和受考核的家庭身上。

\begin{sourcenote}[家族上升的材料，能看见什么]
《春秋》为鲁昭公出奔、阳虎出奔与陈恒弑君保存了简明的年代骨架；《左传》把城邑争夺、人物言辞和道德判断展开为连续叙事，不能当作逐场会议的记录。晋阳与三家分晋的完整故事主要见于《史记》《战国策》及后出的综合编年，田氏代齐又可参照《史记·田敬仲完世家》与周廷册命传统。吴起在楚的措施和遇害，最系统的传世说法出自《史记·孙子吴起列传》，具体条文、先后和动机仍有讨论。因而本节只把它们作为资源、依附与名分如何相互牵动的比较线索，不把任何一家家谱、谋议或制度清单补写成确证事实。
\end{sourcenote}

\begin{turningpoint}[制度得失：让资源脱离家门，还是让国家更深进入家门]
\term{得}：春秋的卿族、封邑和家臣网络能在交通、文书与常备官僚都有限的条件下组织征发和地方治理；战国的重组则试图降低国君对单一家门的依赖。

\term{失}：前一种安排会使军政资源私家化，后一种安排则可能使税、役、刑与战争更直接地抵达普通家庭。齐的易姓、晋的分裂、鲁的空心化和楚的抵抗都说明，国家转型不是简单的“削弱贵族”，而是重新决定谁承担动员成本、谁能把资源说成公共事务。

\term{后续压力}：战国国家会继续竞争谁能更稳定地登记人口、征取赋役、任用官吏并维持军队；它们也将面对一个更尖锐的问题：当国家能力集中之后，怎样防止新的权力中心再次把公共资源据为己有。
\end{turningpoint}
